\documentclass[aspectratio=169]{beamer}

% ==== Theme & basic look ====
\usetheme{Madrid}
\usecolortheme{seahorse}
\setbeamertemplate{navigation symbols}{}
\setbeamertemplate{footline}[frame number]


% Packages for mathematics  
\usepackage{amsmath}  
\usepackage{amsfonts}  
\usepackage{amssymb}  
\usepackage{CJKutf8}
% Support for \Tilde macro used throughout
\providecommand{\Tilde}[1]{\tilde{#1}}
\providecommand{\proofname}{Proof}
\newenvironment<>{proofs}[1][\proofname]{%
    \par
    \def\insertproofname{#1.}%
    \usebeamertemplate{proof begin}#2}
  {\usebeamertemplate{proof end}}
\makeatother

% Math shorthand and vector symbols
\newcommand{\grad}{\nabla}
\newcommand{\EE}{\mathbb{E}}
\newcommand{\dd}{\mathrm{d}}
\newcommand{\kbT}{k_B T}
\newcommand{\vx}{\mathbf{x}}

% Bold vector/time-indexed variables for autoregressive models
\newcommand{\vxt}{\mathbf{x}_t}
\newcommand{\vxone}{\mathbf{x}_1}
\newcommand{\vxtmone}{\mathbf{x}_{t-1}}
\newcommand{\vxtmoneone}{\mathbf{x}_{t-1:1}}
\newcommand{\vxoneT}{\mathbf{x}_{1:T}}
\newcommand{\vtheta}{\boldsymbol{\theta}}
\newcommand{\vphi}{\boldsymbol{\phi}}
\newcommand{\vz}{\mathbf{z}}
\newcommand{\vc}{\mathbf{c}}

% Title page details:   
\title[Understanding autoregressive models]{Understanding autoregressive models from a mathematical view}  
\author{Pipi Hu}  
\institute[BIMSA]{Beijing Institute of Mathematical Sciences and Applications}  
\date{\today}  

\pgfdeclareimage[width=2cm]{logo}{figs/logo.png}
\logo{\pgfuseimage{logo}}

% Outline slide (moved to later; avoid frame before \begin{document})

\setbeamertemplate{navigation symbols}{}
\AtBeginSection[]  
{  
    \begin{frame}<beamer>{Outline}         
        \tableofcontents[currentsection]  
    \end{frame}  
}  
  
\AtBeginSubsection[]  
{  
    \begin{frame}<beamer>{Outline}         
        \tableofcontents[currentsection,currentsubsection]  
    \end{frame}  
}  

\begin{document}  
  
% Title slide  
\begin{frame}  
  \titlepage  
\end{frame}  

\begin{frame}{Feynman's words}
\centering
    What I cannot create, I do not understand.
\end{frame}  
  
% Presentation slides  

\section{Preliminary: Autoregressive Models and Sequential Generation}

\begin{frame}{What are autoregressive models?}
Given a sequence $\vxoneT = (x_1, x_2, \ldots, x_T)$, autoregressive models factorize the joint probability as:
\begin{equation}
    p(\vxoneT) = \prod_{t=1}^T p(x_t | \vxtmoneone)
\end{equation}
where $\vxtmoneone = (x_{t-1}, x_{t-2}, \ldots, x_1)$ represents the history.

Key properties:
\begin{itemize}
    \item \textcolor{red}{Sequential generation}: Generate one token at a time
    \item \textcolor{red}{Conditional independence}: Each token depends only on previous tokens
    \item \textcolor{red}{Causal structure}: No future information leakage
\end{itemize}
\end{frame}

\begin{frame}{Mathematical foundation: Chain rule of probability}
The autoregressive factorization follows directly from the chain rule of probability:
\begin{equation}
    p(\vxoneT) = p(x_1) \prod_{t=2}^T p(x_t | \vxtmoneone)
\end{equation}

This factorization is:
\begin{itemize}
    \item \textbf{Always valid}: No assumptions required
    \item \textbf{Exact}: No approximation involved
    \item \textbf{Order-dependent}: Different orderings give different factorizations
\end{itemize}

The challenge is modeling the conditional distributions $p(x_t | \vxtmoneone)$.
\end{frame}

\begin{frame}{Why autoregressive models?}
Advantages of autoregressive modeling:
\begin{enumerate}
    \item \textbf{Simplicity}: Direct application of maximum likelihood
    \item \textbf{Interpretability}: Clear probabilistic interpretation
    \item \textbf{Flexibility}: Can model complex dependencies
    \item \textbf{Training stability}: Well-understood optimization landscape
\end{enumerate}

Challenges:
\begin{enumerate}
    \item \textbf{Sequential generation}: Cannot parallelize during inference
    \item \textbf{Error propagation}: Mistakes compound over time
    \item \textbf{Long-range dependencies}: Hard to capture distant relationships
\end{enumerate}
\end{frame}

\section{Training Autoregressive Models: Maximum Likelihood}

\begin{frame}{Maximum Likelihood Estimation}
Given a dataset $\mathcal{D} = \{\vxoneT^{(i)}\}_{i=1}^N$, we maximize the log-likelihood:
\begin{equation}
    \mathcal{L}(\vtheta) = \sum_{i=1}^N \log p(\vxoneT^{(i)}; \vtheta)
\end{equation}

Substituting the autoregressive factorization:
\begin{equation}
    \mathcal{L}(\vtheta) = \sum_{i=1}^N \sum_{t=1}^T \log p(x_t^{(i)} | \vxtmoneone^{(i)}; \vtheta)
\end{equation}

This decomposes into independent optimization problems for each conditional distribution.
\end{frame}

\begin{frame}{Cross-entropy loss and its interpretation}
For discrete sequences, the loss becomes:
\begin{equation}
    \mathcal{L}(\vtheta) = -\sum_{i=1}^N \sum_{t=1}^T \log p(x_t^{(i)} | \vxtmoneone^{(i)}; \vtheta)
\end{equation}

This is equivalent to minimizing the cross-entropy:
\begin{equation}
    H(p_{\text{data}}, p_{\vtheta}) = \EE_{\vxoneT \sim p_{\text{data}}} \left[ -\sum_{t=1}^T \log p(x_t | \vxtmoneone; \vtheta) \right]
\end{equation}

\textbf{Key insight}: We're learning to predict the next token given the history, which is exactly what we need for generation.
\end{frame}

\begin{frame}{Teacher forcing and training efficiency}
During training, we use \textbf{teacher forcing}:
\begin{itemize}
    \item Input: Ground truth history $\vxtmoneone$
    \item Target: Next token $x_t$
    \item Loss: Cross-entropy between predicted and true distributions
\end{itemize}

This allows:
\begin{enumerate}
    \item \textbf{Parallel training}: All time steps computed simultaneously
    \item \textbf{Stable gradients}: No error propagation during training
    \item \textbf{Efficient optimization}: Standard backpropagation
\end{enumerate}

\textcolor{red}{Note}: Training and inference use different procedures!
\end{frame}

\section{Architectures for Autoregressive Models}

\begin{frame}{Recurrent Neural Networks (RNNs)}
The simplest autoregressive architecture:
\begin{equation}
    h_t = f(h_{t-1}, x_{t-1}; \vtheta)
\end{equation}
\begin{equation}
    p(x_t | \vxtmoneone) = g(h_t; \vtheta)
\end{equation}

where $h_t$ is the hidden state and $f, g$ are neural networks.

\textbf{Advantages}:
\begin{itemize}
    \item Natural sequential processing
    \item Parameter sharing across time
    \item Variable-length sequences
\end{itemize}

\textbf{Limitations}:
\begin{itemize}
    \item Vanishing/exploding gradients
    \item Limited long-range dependencies
    \item Sequential computation
\end{itemize}
\end{frame}

\begin{frame}{Long Short-Term Memory (LSTM)}
LSTM addresses RNN limitations with gating mechanisms:
\begin{align}
    f_t &= \sigma(W_f \cdot [h_{t-1}, x_t] + b_f) \quad \text{(forget gate)} \\
    i_t &= \sigma(W_i \cdot [h_{t-1}, x_t] + b_i) \quad \text{(input gate)} \\
    \tilde{C}_t &= \tanh(W_C \cdot [h_{t-1}, x_t] + b_C) \quad \text{(candidate values)} \\
    C_t &= f_t * C_{t-1} + i_t * \tilde{C}_t \quad \text{(cell state)} \\
    o_t &= \sigma(W_o \cdot [h_{t-1}, x_t] + b_o) \quad \text{(output gate)} \\
    h_t &= o_t * \tanh(C_t) \quad \text{(hidden state)}
\end{align}

\textbf{Key innovations}:
\begin{itemize}
    \item \textcolor{red}{Gating}: Control information flow
    \item \textcolor{red}{Cell state}: Long-term memory storage
    \item \textcolor{red}{Gradient flow}: Better backpropagation
\end{itemize}
\end{frame}

\begin{frame}{Transformer: Attention-based autoregressive models}
Transformers use self-attention instead of recurrence:
\begin{equation}
    \text{Attention}(Q, K, V) = \text{softmax}\left(\frac{QK^T}{\sqrt{d_k}}\right)V
\end{equation}

For autoregressive generation, we use \textbf{masked attention}:
\begin{equation}
    A_{ij} = \begin{cases}
        \frac{\exp(q_i^T k_j / \sqrt{d_k})}{\sum_{l=1}^i \exp(q_i^T k_l / \sqrt{d_k})} & \text{if } j \leq i \\
        0 & \text{if } j > i
    \end{cases}
\end{equation}

\textbf{Advantages}:
\begin{itemize}
    \item \textcolor{red}{Parallel training}: All positions computed simultaneously
    \item \textcolor{red}{Long-range dependencies}: Direct attention connections
    \item \textcolor{red}{Scalability}: Efficient with modern hardware
\end{itemize}
\end{frame}

\section{Sampling and Generation Strategies}

\begin{frame}{Greedy decoding}
The simplest generation strategy:
\begin{equation}
    x_t = \arg\max_{x} p(x | \vxtmoneone; \vtheta)
\end{equation}

\textbf{Properties}:
\begin{itemize}
    \item Deterministic: Same input always produces same output
    \item Fast: Single forward pass per token
    \item Often suboptimal: Local decisions may not be globally optimal
\end{itemize}

\textbf{Issues}:
\begin{itemize}
    \item \textcolor{red}{Repetition}: May get stuck in loops
    \item \textcolor{red}{Lack of diversity}: Always chooses most likely token
    \item \textcolor{red}{Mode collapse}: May not explore the full distribution
\end{itemize}
\end{frame}

\begin{frame}{Random sampling}
Sample from the full distribution:
\begin{equation}
    x_t \sim p(x | \vxtmoneone; \vtheta)
\end{equation}

\textbf{Advantages}:
\begin{itemize}
    \item \textcolor{red}{Diversity}: Explores the full distribution
    \item \textcolor{red}{Natural}: Reflects model uncertainty
    \item \textcolor{red}{Flexibility}: Can control randomness
\end{itemize}

\textbf{Challenges}:
\begin{itemize}
    \item \textcolor{red}{Quality control}: May generate low-quality sequences
    \item \textcolor{red}{Consistency}: Less predictable output
    \item \textcolor{red}{Evaluation}: Harder to assess quality
\end{itemize}
\end{frame}

\begin{frame}{Temperature scaling and nucleus sampling}
\textbf{Temperature scaling}:
\begin{equation}
    p_{\text{temp}}(x | \vxtmoneone) = \frac{\exp(f(x, \vxtmoneone) / \tau)}{\sum_{x'} \exp(f(x', \vxtmoneone) / \tau)}
\end{equation}

where $\tau > 0$ controls randomness:
\begin{itemize}
    \item $\tau \to 0$: Greedy decoding
    \item $\tau = 1$: Original distribution
    \item $\tau > 1$: More random
\end{itemize}

\textbf{Nucleus sampling} (Holtzman et al., 2019):
\begin{equation}
    \text{Select smallest } V \text{ such that } \sum_{x \in V} p(x | \vxtmoneone) \geq p
\end{equation}
Then sample from the truncated distribution.
\end{frame}

\section{Theoretical Analysis: Expressiveness and Limitations}

\begin{frame}{Universal approximation for autoregressive models}
\begin{theorem}
For any probability distribution $p(\vxoneT)$ over discrete sequences, there exists an autoregressive model that can approximate it arbitrarily well.
\end{theorem}

\begin{proof}
The autoregressive factorization is exact:
\begin{equation}
    p(\vxoneT) = \prod_{t=1}^T p(x_t | \vxtmoneone)
\end{equation}

If we can approximate each conditional $p(x_t | \vxtmoneone)$ arbitrarily well, then the joint distribution is approximated arbitrarily well. Since neural networks are universal function approximators, the result follows.
\end{proof}
\end{frame}

\begin{frame}{The curse of dimensionality in autoregressive models}
Consider a sequence of length $T$ with vocabulary size $V$. The number of possible sequences is $V^T$.

\textbf{Exponential growth}:
\begin{itemize}
    \item $T = 10, V = 1000$: $10^{30}$ sequences
    \item $T = 100, V = 50000$: $50000^{100}$ sequences
\end{itemize}

\textbf{Implications}:
\begin{enumerate}
    \item \textcolor{red}{Data sparsity}: Most sequences never seen in training
    \item \textcolor{red}{Generalization challenge}: Hard to generalize to unseen patterns
    \item \textcolor{red}{Computational limits}: Cannot store or compute exact probabilities
\end{enumerate}

\textbf{Why autoregressive models work}:
\begin{itemize}
    \item \textcolor{red}{Structure}: Real sequences have structure and patterns
    \item \textcolor{red}{Conditional independence}: Each step depends only on history
    \item \textcolor{red}{Neural approximation}: Networks can learn complex patterns
\end{itemize}
\end{frame}

\begin{frame}{Information-theoretic perspective}
The entropy of an autoregressive model is:
\begin{equation}
    H(p_{\vtheta}) = \EE_{\vxoneT \sim p_{\vtheta}} \left[ -\sum_{t=1}^T \log p(x_t | \vxtmoneone; \vtheta) \right]
\end{equation}

\textbf{Key insights}:
\begin{itemize}
    \item \textcolor{red}{Conditional entropy}: $H(X_t | X_{t-1:1}) \leq H(X_t)$
    \item \textcolor{red}{Chain rule}: $H(X_{1:T}) = \sum_{t=1}^T H(X_t | X_{t-1:1})$
    \item \textcolor{red}{Compression}: Autoregressive models compress information
\end{itemize}

\textbf{Compression ratio}:
\begin{equation}
    \text{Compression} = \frac{H(X_{1:T})}{T \log V} = \frac{\sum_{t=1}^T H(X_t | X_{t-1:1})}{T \log V}
\end{equation}

Lower compression ratio means better modeling of dependencies.
\end{frame}

\section{Advanced Topics: Variational Autoregressive Models}

\begin{frame}{Variational Autoencoders for sequences}
Standard autoregressive models are deterministic. We can add stochasticity:

\begin{equation}
    p(\vxoneT) = \int p(\vxoneT | \vz) p(\vz) d\vz
\end{equation}

where $\vz$ is a latent variable and:
\begin{equation}
    p(\vxoneT | \vz) = \prod_{t=1}^T p(x_t | \vxtmoneone, \vz; \vtheta)
\end{equation}

\textbf{Training objective} (ELBO):
\begin{equation}
    \mathcal{L} = \EE_{q(\vz | \vxoneT)} \left[ \log p(\vxoneT | \vz) \right] - D_{KL}(q(\vz | \vxoneT) \| p(\vz))
\end{equation}

\textbf{Benefits}:
\begin{itemize}
    \item \textcolor{red}{Controllable generation}: Vary $\vz$ for different styles
    \item \textcolor{red}{Better modeling}: Can capture multiple modes
    \item \textcolor{red}{Interpretability}: Latent space has semantic meaning
\end{itemize}
\end{frame}

\begin{frame}{Conditional autoregressive models}
Instead of unconditional generation, we can condition on additional information:

\begin{equation}
    p(\vxoneT | \vc) = \prod_{t=1}^T p(x_t | \vxtmoneone, \vc; \vtheta)
\end{equation}

where $\vc$ is conditioning information (e.g., class label, prompt, style).

\textbf{Applications}:
\begin{itemize}
    \item \textcolor{red}{Controlled generation}: Style transfer, conditional completion
    \item \textcolor{red}{Multi-modal learning}: Text-to-image, image-to-text
    \item \textcolor{red}{Few-shot learning}: Adapt to new domains quickly
\end{itemize}

\textbf{Training}:
\begin{equation}
    \mathcal{L} = \EE_{(\vxoneT, \vc) \sim \mathcal{D}} \left[ -\sum_{t=1}^T \log p(x_t | \vxtmoneone, \vc; \vtheta) \right]
\end{equation}
\end{frame}

\section{Evaluation and Challenges}

\begin{frame}{Evaluation metrics for autoregressive models}
\textbf{Perplexity}:
\begin{equation}
    \text{Perplexity} = \exp\left(-\frac{1}{T} \sum_{t=1}^T \log p(x_t | \vxtmoneone; \vtheta)\right)
\end{equation}

\textbf{Other metrics}:
\begin{itemize}
    \item \textcolor{red}{BLEU}: N-gram overlap with reference
    \item \textcolor{red}{ROUGE}: Recall-oriented evaluation
    \item \textcolor{red}{BERTScore}: Semantic similarity using BERT
    \item \textcolor{red}{Human evaluation}: Subjective quality assessment
\end{itemize}

\textbf{Challenges}:
\begin{itemize}
    \item \textcolor{red}{No single metric}: Different tasks need different metrics
    \item \textcolor{red}{Reference dependence}: Many metrics need reference text
    \item \textcolor{red}{Human evaluation}: Expensive and subjective
\end{itemize}
\end{frame}

\begin{frame}{Common challenges and limitations}
\textbf{Training challenges}:
\begin{enumerate}
    \item \textcolor{red}{Exposure bias}: Training uses ground truth, inference uses model predictions
    \item \textcolor{red}{Teacher forcing mismatch}: Different distributions during training and inference
    \item \textcolor{red}{Long sequences}: Gradient vanishing/exploding
\end{enumerate}

\textbf{Generation challenges}:
\begin{enumerate}
    \item \textcolor{red}{Repetition}: Models may repeat phrases or sentences
    \item \textcolor{red}{Coherence}: Long-range dependencies hard to maintain
    \item \textcolor{red}{Diversity vs. Quality}: Trade-off between exploration and exploitation
\end{enumerate}

\textbf{Computational challenges}:
\begin{enumerate}
    \item \textcolor{red}{Sequential generation}: Cannot parallelize inference
    \item \textcolor{red}{Memory requirements}: Need to store full history
    \item \textcolor{red}{Scalability}: Training large models is expensive
\end{enumerate}
\end{frame}

\section{Recent Advances and Future Directions}

\begin{frame}{Non-autoregressive models}
\textbf{Motivation}: Autoregressive models are slow due to sequential generation.

\textbf{Non-autoregressive approaches}:
\begin{itemize}
    \item \textcolor{red}{Parallel decoding}: Generate all tokens simultaneously
    \item \textcolor{red}{Iterative refinement}: Multiple passes to improve quality
    \item \textcolor{red}{Latent variable models}: Use latent representations
\end{itemize}

\textbf{Challenges}:
\begin{itemize}
    \item \textcolor{red}{Quality degradation}: Parallel generation often lower quality
    \item \textcolor{red}{Training complexity}: Harder to train effectively
    \item \textcolor{red}{Evaluation}: Different metrics needed
\end{itemize}

\textbf{Examples}:
\begin{itemize}
    \item \textcolor{red}{NAT (Gu et al., 2017)}: Non-autoregressive neural machine translation
    \item \textcolor{red}{GLAT (Qian et al., 2021)}: Glancing transformer for NAT
    \item \textcolor{red}{CMLM (Ghazvininejad et al., 2019)}: Conditional masked language model
\end{itemize}
\end{frame}

\begin{frame}{Large language models and scaling laws}
\textbf{Scaling laws} (Kaplan et al., 2020):
\begin{equation}
    L(N, D) = \left(\frac{N_c}{N}\right)^{\alpha_N} + \left(\frac{D_c}{D}\right)^{\alpha_D}
\end{equation}

where $N$ is model size, $D$ is dataset size, and $L$ is loss.

\textbf{Key findings}:
\begin{itemize}
    \item \textcolor{red}{Power law scaling}: Performance improves predictably with scale
    \item \textcolor{red}{Emergent abilities}: New capabilities emerge at certain scales
    \item \textcolor{red}{Efficient scaling}: Larger models are more sample-efficient
\end{itemize}

\textbf{Implications}:
\begin{itemize}
    \item \textcolor{red}{Compute requirements}: Training costs scale dramatically
    \item \textcolor{red}{Data requirements}: Need massive datasets
    \item \textcolor{red}{Architecture choices}: Transformer scales well
\end{itemize}
\end{frame}

\section{Open Questions and Research Directions}

\begin{frame}{Open questions}
If you are interested in these questions, please feel free to write an email to me {\color{red}(\href{mailto:pisquare@microsoft.com}{pisquare@microsoft.com})} with your comments.

\begin{enumerate}[(1.)]
    \item How can we design better non-autoregressive models that maintain quality while enabling parallel generation?
    \item What are the fundamental limits of autoregressive modeling? Can we prove lower bounds on the required model size for certain tasks?
    \item How do we balance the trade-off between generation speed and quality in practical applications?
    \item What role do autoregressive models play in the broader landscape of generative AI, and how do they compare to diffusion models and other approaches?
    \item Can we develop theoretical frameworks to understand when autoregressive models will succeed or fail?
\end{enumerate}
\end{frame}


% Thank you slide  
\begin{frame}{Welcome inputs}  
  \centering  
  Any comments?
  \\[2em]
  Welcome your inputs!  
\end{frame}  
  
\end{document}
